% Exposé V: back to the volume's parf
\renewcommand{\parf}{\operatorname{parf}}
\sgaStableTarget{sga6.e05.expose}\section*{\underline{Exposé V}\\[3pt] \underline{\underline{GENERALITIES ON @@M0@@-RINGS}}}
\begin{center}
\underline{by P. BERTHELOT}
\end{center}

The theory of @@M1@@-rings is useful whenever one uses rings such as the ring of classes of vector bundles on a ringed site, the representation ring of a group, etc..., for which the exterior-power operations define maps @@M2@@. We shall use it extensively in the remainder of the Seminar, so it is necessary first to develop its general formalism. This exposé is thus entirely independent of the rest of the seminar and can serve users of @@M3@@-rings outside algebraic geometry.

The first section essentially gives Hirzebruch's theory of ``multiplicative sequences of polynomials'' (cf. [2]); the functorial viewpoint adopted here has the advantage, besides its generality, of considerably reducing the calculations. The following sections will show how, to every augmented @@M4@@-ring, one can associate a ``Chern ring'' that permits the development of a theory of Chern classes. This is the theory that will enable us to state the Riemann--Roch theorem in the general form of \hyperlink{sga6.e08.expose}{Exp. VIII}.

\sgaStableTarget{sga6.e05.sec.1}\section*{1. Universal polynomials}

\sgaStableTarget{sga6.e05.para.1-1}1.1. Let @@M5@@ be a commutative ring with unit, fixed throughout this section. We denote by @@M6@@ the category of filtered commutative @@M7@@-algebras with unit, whose filtration is decreasing and such that the group of elements of filtration zero is the whole algebra. We define a functor on @@M8@@, with values in the category of groups
